
\section{MIO Modelos de Transcripción}
\label{section:6-1-Modelos de Transcripción}

En esta sección se detalla nuestra experimentación en cuanto a los métodos de transcripción.

\subsection{Proceso de Evaluación}
\label{subsection:6-1-Proceso de Evaluación}

La evaluación del método de transcripción deberá tener un buen balance en los métodos mencionados en \autoref{subsection:3-RNF3: Evaluación de Herramientas de Transcripción}.

Como ya fue mencionado en la sección, para evaluar el método de transcripción debemos partir de una transcripción perfecta(\textit{ground-truth}). Para esto, se transcribieron manualmente las clases \href{https://open.fing.edu.uy/courses/comp1-2024/2/},{Computación 1 Clase 1-Teórico 2024}
\href{https://open.fing.edu.uy/courses/comp1-2024/1/}{Computación 1 Clase 2-Teórico 2024}, \href{https://open.fing.edu.uy/courses/pln/1}{Introducción Procesamiento de Lenguaje Natural Clase 1, Teórico 2016} y \href{https://open.fing.edu.uy/courses/pln/1/}{Introducción Procesamiento de Lenguaje Natural Clase 2, Teórico 2016}

\subsection{Métricas}
En la siguiente subsección se presentan las métricas encontradas durante la comparativa en la tabla

\begin{table}[htbp]
\centering
\caption{Model Evaluation Metrics (ad = autodetect, es = español)}
\begin{tabular}{l l l S[table-format=1.3] S[table-format=1.3] S[table-format=1.3] S[table-format=1.3] S[table-format=1.3]}
\toprule
\textbf{Model} & \textbf{Test Set} & \textbf{Lang} & \textbf{WER} & \textbf{MER} & \textbf{WIL} & \textbf{WIP} & \textbf{CER} \\
\midrule
medium & clase 1 & ad & 0.046 & 0.045 & 0.071 & 0.928 & 0.023 \\
       &               & es & 0.046 & 0.045 & 0.071 & 0.928 & 0.023 \\
\cmidrule(r){1-8}
medium & clase 2 & ad & 0.086 & 0.085 & 0.120 & 0.879 & 0.052 \\
       &               & es & 0.042 & 0.042 & 0.057 & 0.942 & 0.027 \\
\midrule
small  & clase 1 & ad & 0.150 & 0.148 & 0.248 & 0.751 & 0.054 \\
       &               & es & 0.150 & 0.148 & 0.248 & 0.751 & 0.054 \\
\cmidrule(r){1-8}
small  & clase 2 & ad & 0.139 & 0.137 & 0.219 & 0.780 & 0.059 \\
       &               & es & 0.234 & 0.228 & 0.372 & 0.627 & 0.081 \\
\midrule
base   & clase 1 & ad & 0.203 & 0.199 & 0.324 & 0.675 & 0.076 \\
       &               & es & 0.198 & 0.194 & 0.317 & 0.682 & 0.074 \\
\bottomrule
\end{tabular}
\
\end{table}
Añadir análisis de la tabla.
\subsection{Análisis Cualitativo}
\label{subsection:6-1-análisis Cualitativo}

Durante la comparativa entre la transcripción \textit{ground-truth} y la comparativa real si bien se encuentran errores, varios errores son localizados y frecuentes en situaciones especificas que no afectan significativamente la .
Las muletillas (\textit{filler words}) en ingles, aquellas palabras o sonidos utilizados mientras pensamos lo siguiente que diremos, no suelen ser captadas precisamente por el ASR. 

Es cierto que la mejor ASR debería de poder captar \textit{filler-words} exitosamente, sin embargo, como en el proyecto se utilizan dichas transcripciones para un sistema RAG, se puede argumentar que terminar que los texto quitandoles muletillas son mejor, debido a que no se pierde contexto en palabras que no aportan al significado de lo dicho. Si bien no ha sido realizado, es relevante poder presentar (referencia a tabla anterior) quitandole las muletillas de lo mencionado. 
\subsection{Resultados}
\label{subsection:6-1-Resultados}
Los resultados entra

\begin{table}[!ht]
\centering
\caption{Evaluación de Modelos }
\begin{tabular}{l l l S[table-format=1.3] S[table-format=1.3] S[table-format=1.3] S[table-format=1.3] S[table-format=1.3]}
\toprule
\textbf{Model} & \textbf{Test Set} & \textbf{Lang} & \textbf{WER} & \textbf{MER} & \textbf{WIL} & \textbf{WIP} & \textbf{CER} \\
\midrule
medium & clase 1 & ad & 0.046 & 0.045 & 0.071 & 0.928 & 0.023 \\
       &               & es & 0.046 & 0.045 & 0.071 & 0.928 & 0.023 \\
\cmidrule(r){1-8}
medium & clase 2 & ad & 0.086 & 0.085 & 0.120 & 0.879 & 0.052 \\
       &               & es & 0.042 & 0.042 & 0.057 & 0.942 & 0.027 \\
\midrule
small  & clase 1 & ad & 0.150 & 0.148 & 0.248 & 0.751 & 0.054 \\
       &               & es & 0.150 & 0.148 & 0.248 & 0.751 & 0.054 \\
\cmidrule(r){1-8}
small  & clase 2 & ad & 0.139 & 0.137 & 0.219 & 0.780 & 0.059 \\
       &               & es & 0.234 & 0.228 & 0.372 & 0.627 & 0.081 \\
\midrule
base   & clase 1 & ad & 0.203 & 0.199 & 0.324 & 0.675 & 0.076 \\
       &               & es & 0.198 & 0.194 & 0.317 & 0.682 & 0.074 \\
\bottomrule
\end{tabular}
\end{table}


\section{Expansión de la Base de Datos}
\label{section:6-2-Expansión de la Base de Datos}
\todo[inline,color=cyan]{Añadir todo el Volumen de \Gls{OpenFing}}


\section{\textit{Question Answering}}
\label{section:6-2-Question Answering}

\subsection{Proceso de Evaluación}
Para evaluar la respuesta de cada uno de los modelos, se formularon preguntas correspondientes a cuatro categorías diferenciadas:

\begin{itemize}
    \item \textbf{Presente}: La respuesta se encuentra explícitamente en el texto.Se pueden 
    \item \textbf{Derivable}: La respuesta no aparece de forma explícita, pero resulta deducible a partir del \gls{contexto} proporcionado.
    \item \textbf{No presente} ni derivable: La respuesta no se encuentra explícitamente ni es deducible del \gls{contexto} suministrado.
    \item \textbf{Irrelacionada}: La respuesta carece de relación con la temática del \gls{contexto} proporcionado.
\end{itemize}

Por ejemplo estas son algunas de las preguntas de cada una de las categorias que propusimos pata la Clase 1 de Procesamiento Natural de Lenguaje.

